\documentclass[11pt]{article}
\usepackage{../common/dastyle}
\graphicspath{{../figures/}}

\begin{document}
\pagestyle{fancy}
\sheethead{TD -- Adaptation de Domaine}{Ce TD reprend et prolonge les notions du cours (séance 6) : repondération d'instances sous covariate shift, estimation du ratio de densités (KDE, KMM, classifieur de domaine), puis bornes de généralisation ($\Hcal\Delta\Hcal$, PAC-Bayes) et un aperçu des méthodes profondes récentes (DANN). Les exercices marqués \textbf{(Application)} demandent un calcul numérique court ; ceux marqués \textbf{(Preuve)} une démonstration. La correction est fournie dans un document séparé (\texttt{td\_da\_solutions.pdf}).}

\section{Fondamentaux : covariate shift et repondération}

\begin{exercise}[Identifier l'hypothèse de correspondance]
Pour chacune des situations suivantes, indiquez laquelle des quatre hypothèses vues en cours
(covariate shift, target shift, biais de sélection, concept shift) est la plus plausible, et
justifiez en une phrase.
\begin{enumerate}[label=(\alph*)]
\item Un hôpital A forme un modèle de dépistage sur ses patients (population plutôt âgée) ; l'hôpital B,
dans une autre région, a une population plus jeune, mais la même pathologie se manifeste de la même façon
cliniquement dans les deux populations.
\item Une plateforme de recommandation de films est entraînée sur des clics de 2015, où les
utilisateurs préféraient les DVD physiques loués en magasin ; en 2025, les mêmes types de contenus sont
préférés, mais la manière de les consommer (SVOD) a totalement changé la notion même de \enquote{clic}.
\item Un capteur de température, mal calibré, sous-échantillonne systématiquement les mesures extrêmes
lors de la collecte des données d'entraînement, alors que le capteur de production (bien calibré) ne
souffre pas de ce biais.
\end{enumerate}
\end{exercise}

\begin{exercise}[Régression pondérée -- calcul explicite] \textbf{(Application)}
\label{ex:td-weighted-reg}
On considère $P_S(X) = \mathcal N(-1,1)$, $P_T(X) = \mathcal N(1, 1)$, et le modèle
$P_S(Y\mid X=x) = P_T(Y\mid X=x) = \mathcal N(|x|, \sigma^2)$ (covariate shift).
\begin{enumerate}[label=(\alph*)]
\item Rappelez la définition de $w(x) = p_t(x)/p_s(x)$ et montrez que, dans ce cas précis,
$w(x) = \exp(2x)$.
\item On observe un unique point source $x_1 = -0.5$. Calculez $w(x_1)$.
\item On observe un second point source $x_2 = 1.5$. Calculez $w(x_2)$ et comparez à $w(x_1)$. Que
constatez-vous sur le comportement de $w$ loin de la zone où vivent les données source ?
\item \textbf{(Preuve)} Pour $X\in\R^{m\times p}$, $Y\in\R^m$, et $\Delta=\mathrm{diag}(\sqrt{w_1},
\ldots,\sqrt{w_m})$, retrouvez la solution $\beta^\star$ du problème de régression pondérée
$\min_\beta \sum_i w_i(x_i\beta^\top - y_i)^2$.
\end{enumerate}
\end{exercise}

\begin{exercise}[Effective Sample Size] \textbf{(Application)}
\label{ex:ess}
Un problème pratique des méthodes instance-based : si quelques poids $w_i$ dominent tous les autres, on
n'utilise \enquote{effectivement} qu'une poignée de points source, même si $m$ est grand. On mesure cela
par la taille d'échantillon effective :
\[
\mathrm{ESS}(w) = \frac{\left(\sum_{i=1}^m w_i\right)^2}{\sum_{i=1}^m w_i^2}.
\]
\begin{enumerate}[label=(\alph*)]
\item Calculez $\mathrm{ESS}(w)$ pour $m=4$ poids égaux $w=(1,1,1,1)$. Commentez.
\item Calculez $\mathrm{ESS}(w)$ pour $w=(10, 0.1, 0.1, 0.1)$. Commentez, en particulier au regard de la
Figure~3.5 du cours (supports peu recouvrants).
\item \textbf{(Preuve)} Montrez que $1 \le \mathrm{ESS}(w) \le m$ pour tout vecteur de poids positifs
$w$ (indication : utilisez l'inégalité de Cauchy-Schwarz pour la borne supérieure).
\end{enumerate}
\end{exercise}

\section{Estimer le ratio de densités}

\begin{exercise}[Densité paramétrique vs. noyaux] \textbf{(Application)}
On dispose de $S_X = \{-1.2, -0.9, -1.1, -0.8\}$ et $T_X=\{0.8, 1.3, 1.0\}$.
\begin{enumerate}[label=(\alph*)]
\item Calculez les estimateurs de moyenne et variance empiriques $(\hat\mu_s,\hat\sigma_s^2)$ et
$(\hat\mu_t,\hat\sigma_t^2)$ utilisés par l'approche paramétrique gaussienne (méthode du maximum de
vraisemblance).
\item En déduire $\hat w(x_i) = \hat p_t(x_i)/\hat p_s(x_i)$ pour $x_1=-1.2$.
\item \textbf{(Preuve)} Pour un noyau gaussien $k(x,x')=\exp(-(x-x')^2/2\sigma^2)$, écrivez l'estimateur
de densité par noyaux $\hat p_s(x)$ non normalisé pour l'échantillon $S_X$ ci-dessus, et donnez sa valeur
en $x=-1$ pour $\sigma=0.3$ (à $10^{-2}$ près).
\end{enumerate}
\end{exercise}

\begin{exercise}[Kernel Mean Matching] \textbf{(Preuve)}
\label{ex:kmm}
On rappelle que KMM résout $w^\star = \argmin_{w\ge0,\ w^\top\mathbf 1 = m}\ w^\top K w - 2\kappa^\top w$,
où $K_{ij}=k(x_i,x_j)$, $\kappa_i = \frac{m}{n}\sum_{j} k(x_i,x'_j)$.
\begin{enumerate}[label=(\alph*)]
\item Sans la contrainte $w\ge 0$ ni $w^\top\mathbf1=m$, écrivez la condition d'optimalité du premier
ordre du problème $\min_w w^\top K w - 2\kappa^\top w$ et donnez la solution non contrainte $w^\star$
(en supposant $K$ inversible).
\item Que représente $\kappa_i$ intuitivement ? Pourquoi la normalisation par $m/n$ est-elle nécessaire ?
\item Pourquoi le problème original impose-t-il $w\ge 0$ et $\sum_i w_i = m$ ? Que se passerait-il pour
l'interprétation de $w$ comme ratio de densités si l'on retirait ces contraintes ?
\end{enumerate}
\end{exercise}

\begin{exercise}[Classifieur de domaine] \textbf{(Preuve)}
\label{ex:iwc-td}
On adopte ici la convention \emph{inverse} de celle du cours : $\phi(x) = \widehat\P(\text{domaine}=S
\mid X=x)$ (et non $T$). On entraîne $\phi$ par minimisation de l'entropie croisée
\[
\phi^\star = \argmax_{\phi}\ \Esp{P_S(X)}{\log\phi(X)} + \Esp{P_T(X)}{\log(1-\phi(X))}.
\]
\begin{enumerate}[label=(\alph*)]
\item Montrez que $\phi^\star(x) = p_s(x)/(p_s(x)+p_t(x))$.
\item En déduire l'expression de $w(x)=p_t(x)/p_s(x)$ en fonction de $\phi^\star(x)$ uniquement.
\item Un data scientist utilise par erreur un classifieur de domaine avec la convention du cours
($\phi(x)=\widehat\P(T\mid x)$) mais applique la formule $w(x)=1/\phi(x)-1$ de cet exercice. Quelle erreur
commet-il, et quel est l'effet pratique sur les poids obtenus (sous-estimation ou surestimation de
l'importance des zones \enquote{typiquement cible}) ?
\end{enumerate}
\end{exercise}

\section{Bornes de généralisation}

\begin{exercise}[$\Hcal$-divergence sur un exemple jouet] \textbf{(Application)}
On considère $\X=\{1,2,3,4\}$, avec $D_S$ uniforme sur $\{1,2\}$ et $D_T$ uniforme sur $\{3,4\}$. On prend
$\Hcal = \{h_1, h_2\}$ où $h_1(x) = \1[x\le 2]$ et $h_2(x)=\1[x\le 3]$.
\begin{enumerate}[label=(\alph*)]
\item Calculez $\P_{D_S}[h_1(x)=1]$ et $\P_{D_T}[h_1(x)=1]$.
\item En déduire une valeur de $d_\Hcal(D_S,D_T)$ obtenue avec $h_1$ (rappel :
$d_\Hcal(D_S,D_T)=2\sup_{h\in\Hcal}|\P_{D_S}[h=1]-\P_{D_T}[h=1]|$).
\item Refaites le calcul avec $h_2$. Quel est le classifieur de $\Hcal$ qui \enquote{sépare le mieux}
source et cible ici, et quelle est la valeur de $d_\Hcal(D_S,D_T)$ pour cette classe $\Hcal=\{h_1,h_2\}$ ?
\end{enumerate}
\end{exercise}

\begin{exercise}[Lecture de la borne de Ben-David] \textbf{(Application)}
On suppose qu'un data scientist entraîne un classifieur $h$ et mesure $R_S(h) = 0.05$ (risque source),
un classifieur de domaine avec un taux d'erreur $\mathrm{err}(\phi)=0.42$ pour séparer $S_X$ de $T_X$, et
estime (par une méthode indépendante) $\lambda_\Hcal \approx 0.03$.
\begin{enumerate}[label=(\alph*)]
\item Rappelez la formule reliant $\mathrm{err}(\phi)$ (erreur du meilleur classifieur de domaine) et le
proxy empirique $\hat d_\Hcal(S_X,T_X)$.
\item Calculez la valeur numérique de $\hat d_\Hcal(S_X,T_X)$ pour cet exemple.
\item En déduire une borne numérique sur $R_T(h)$ à l'aide du théorème de Ben-David et al.
\item Un second data scientist obtient un classifieur de domaine bien plus performant,
$\mathrm{err}(\phi')=0.08$, sur les mêmes données. Que peut-on en conclure sur la difficulté du problème
de transfert dans ce second cas ? La borne obtenue est-elle plus ou moins informative ?
\end{enumerate}
\end{exercise}

\begin{exercise}[Comprendre le rôle du $\KL$ dans la borne PAC-Bayésienne] \textbf{(Application)}
On propose deux posteriors $Q_1$ et $Q_2$ sur une même classe $\Hcal$, avec les valeurs suivantes (unités
arbitraires, $m=500$, $\delta=0.05$) :
\begin{center}
\begin{tabular}{lccc}
\toprule
& $R_S(Q)$ & $\mathrm{dis}_T(Q)/2$ & $\KL(Q\|P)$ \\
\midrule
$Q_1$ & $0.10$ & $0.04$ & $2.0$ \\
$Q_2$ & $0.06$ & $0.03$ & $45.0$ \\
\bottomrule
\end{tabular}
\end{center}
On néglige ici le terme $\lambda(Q)$ (supposé identique et faible pour les deux) et on prend le terme de
complexité sous la forme simplifiée $\sqrt{(\KL(Q\|P)+\log(1/\delta))/m}$.
\begin{enumerate}[label=(\alph*)]
\item Calculez la borne PAC-Bayésienne approchée pour $Q_1$ puis pour $Q_2$.
\item Quel posterior semble préférable au vu du risque source seul ? Et au vu de la borne complète ?
\item En une phrase, expliquez pourquoi ce phénomène illustre le compromis biais/complexité déjà rencontré
dans d'autres contextes d'apprentissage statistique.
\end{enumerate}
\end{exercise}

\section{Ouverture : méthodes récentes}

\begin{exercise}[DANN et Gradient Reversal Layer] \textbf{(Preuve/Réflexion)}
\begin{enumerate}[label=(\alph*)]
\item DANN entraîne conjointement un extracteur de features $G_f$, un classifieur de labels $G_y$ et un
classifieur de domaine $G_d$. Écrivez (sans détailler l'architecture) la fonction objectif combinée que
$G_f$ et $G_y$ cherchent à \emph{minimiser}, et celle que $G_d$ cherche à minimiser, en notant
$\mathcal L_y$ la perte de classification des labels et $\mathcal L_d$ la perte du classifieur de domaine.
\item Le Gradient Reversal Layer multiplie le gradient de $\mathcal L_d$ par $-\lambda$ (avec
$\lambda > 0$) lorsqu'il traverse $G_f$ dans la rétropropagation. Expliquez pourquoi cela permet
d'entraîner $G_f$ à \emph{maximiser} $\mathcal L_d$ tout en gardant une seule passe de descente de
gradient pour l'ensemble du réseau.
\item Reliez cette procédure à la borne de Ben-David vue en cours : quelle quantité de cette borne DANN
cherche-t-il concrètement à minimiser via $G_d$ ?
\item \textbf{(Réflexion)} D'après la mise en garde de Zhao et al. (2019) vue en cours, dans quelle
situation (en termes de proportions de classes) le principe même de DANN peut-il devenir contre-productif ?
\end{enumerate}
\end{exercise}

\end{document}
